\documentclass[11pt]{article}
\usepackage[margin=1in]{geometry}
\usepackage{graphicx}
\usepackage{booktabs}

\title{Analisis de Interpolacion en Rotacion y Escalado}
\author{}
\date{2026-03-10}

\begin{document}
\maketitle

\section*{Objetivo}
Evaluar como los metodos de interpolacion (nearest neighbor, bilinear y bicubic)
afectan la calidad de la imagen al rotar y escalar, con enfasis en la legibilidad
de texto.

\section*{Datos y salidas}
Se trabajaron imagenes ubicadas en \texttt{interpolation/imagen}. Los resultados
se guardaron en:
\begin{itemize}
  \item \texttt{manipulated/interpolation}: escalado 2x con cada interpolacion.
  \item \texttt{manipulated/rotations}: rotacion 45 y 90 grados (sin escalar).
  \item \texttt{manipulated/RnS}: rotacion (45, 90) y luego escalado 2x.
  \item \texttt{manipulated/RnSText}: \texttt{text.png} rotado 90 grados y luego
        escalado 2x con cada interpolacion.
  \item \texttt{manipulated/TextInterpolate50}: \texttt{text.png} escalado 50\%.
  \item \texttt{manipulated/TextInterpolate200}: \texttt{text.png} escalado 200\%.
\end{itemize}

\section*{Metodologia}
Se compararon tres metodos de interpolacion:
\begin{itemize}
  \item \textbf{Nearest neighbor}: replica el pixel mas cercano (rapido, pero con
        aliasing).
  \item \textbf{Bilinear}: promedio lineal 2D (suaviza, reduce aliasing).
  \item \textbf{Bicubic}: interpolacion cubica (mejor preservacion de detalle,
        posible halo/ringing).
\end{itemize}

\section*{Resultados y analisis}

\subsection*{Rotacion (imagenes generales)}
\textbf{Como afectan los metodos de interpolacion al resultado de la rotacion?}

En rotaciones no multiples de 90 grados, la interpolacion es critica porque se
re-muestrea la imagen. En general:
\begin{itemize}
  \item \textbf{Nearest} produce bordes dentados (aliasing) en diagonales y
        contornos, especialmente visibles en lineas finas.
  \item \textbf{Bilinear} reduce el aliasing pero suaviza texturas y bordes.
  \item \textbf{Bicubic} mantiene mas detalle que bilinear, aunque puede generar
        ligeros halos alrededor de bordes de alto contraste.
\end{itemize}

\textbf{Que metodo mantiene mejor la calidad de la imagen original?}
Para imagenes naturales, \textbf{bicubic} suele ofrecer la mejor conservacion
de detalle percibido con menos degradacion global. \textbf{Bilinear} es un buen
compromiso entre suavizado y nitidez. \textbf{Nearest} es el menos fiel por el
aliasing visible.

\subsection*{Texto (\texttt{text.png})}
\textbf{Que metodo conserva mejor la legibilidad del texto?}

Para texto de alto contraste:
\begin{itemize}
  \item \textbf{Nearest} mantiene bordes muy definidos pero introduce escalones
        en diagonales y curvas; aun asi, suele conservar la legibilidad porque
        preserva el contraste.
  \item \textbf{Bilinear} suaviza bordes y puede adelgazar trazos finos, lo cual
        reduce ligeramente la legibilidad.
  \item \textbf{Bicubic} puede introducir halos suaves alrededor de letras,
        especialmente al rotar 90 y escalar 2x, lo que puede afectar contornos.
\end{itemize}

En este contexto, \textbf{nearest} suele mantener la legibilidad mejor en texto
pequeno y de alto contraste, aunque con aliasing visible. Si se prioriza un
aspecto mas suave y menos ``pixelado'', \textbf{bilinear} es una alternativa
equilibrada.

\subsection*{Escalado 50\% y 200\% (sin rotacion)}
\textbf{Calidad de la imagen final.}

\begin{itemize}
  \item \textbf{50\%}: bilinear y bicubic reducen aliasing, pero pierden algo de
        nitidez. Nearest mantiene nitidez de bordes a costa de legibilidad.
  \item \textbf{200\%}: bicubic tiende a preservar mas detalle y suavidad, mientras
        bilinear es mas suave pero con menos detalle. Nearest amplifica pixeles y
        se ve mas ``cuadriculado''.
\end{itemize}

\section*{Resumen comparativo}
\begin{center}
\begin{tabular}{@{}llll@{}}
\toprule
Caso & Nearest & Bilinear & Bicubic \\ \midrule
Rotacion & Aliasing visible & Suavizado moderado & Mejor detalle, posible halo \\
Texto (legibilidad) & Alta, con escalones & Media, mas suave & Media, posible halo \\
Escalado 50\% & Aliasing & Suave & Mejor equilibrio \\
Escalado 200\% & Pixelado & Suave & Mejor detalle \\ \bottomrule
\end{tabular}
\end{center}

\section*{Conclusiones}
\begin{itemize}
  \item Para imagenes naturales y rotaciones, \textbf{bicubic} suele conservar
        mejor la calidad percibida.
  \item Para texto, \textbf{nearest} tiende a preservar la legibilidad por el
        contraste y la nitidez de bordes, aunque con aliasing.
  \item \textbf{Bilinear} es un compromiso cuando se busca reducir artefactos a
        costa de un leve suavizado.
\end{itemize}

\end{document}
